\chapter{1988年全国卷（文）}

\section{单选题}

\begin{problem}
\(\left(\frac{1-\i}{1+\i}\right)^2\) 的值等于（\quad）

\choices
    {\(1\)}
    {\(-1\)}
    {\(\i\)}
    {\(-\i\)}
\end{problem}

\begin{answer}
B
\end{answer}

\begin{solution}
因为 \((1+\i)(-\i)=1-\i\)，所以
\(\dfrac{1-\i}{1+\i}=-\i\). 于是
\(\left(\dfrac{1-\i}{1+\i}\right)^2=(-\i)^2=-1\)，故选 B.
\end{solution}


\begin{problem}
圆 \(M\) 的方程为 \(\left(x-3\right)^2 + \left(y-2\right)^2 = 2\)，直线 \(L\) 的方程为 \(x+y-3 = 0\)，点 \(P\) 的坐标为 \((2,1)\)，那么（\quad）

\choices
    {点 \(P\) 在直线 \(L\) 上，但不在圆 \(M\) 上}
    {点 \(P\) 在圆 \(M\) 上，但不在直线 \(L\) 上}
    {点 \(P\) 既在圆 \(M\) 上，又在直线 \(L\) 上}
    {点 \(P\) 既不在直线 \(L\) 上，也不在圆 \(M\) 上}
\end{problem}

\begin{answer}
C
\end{answer}

\begin{solution}
把 \(P\left(2,1\right)\) 代入直线 \(L\) 的方程，得到
\(2+1-3=0\)，所以 \(P\in L\). 再把 \(P\) 的坐标代入圆 \(M\) 的方程左端，有
\(\left(2-3\right)^2+\left(1-2\right)^2=1+1=2\)，等于方程右端，因此 \(P\in M\). 所以 \(P\) 既在圆 \(M\) 上，又在直线 \(L\) 上，故选 C.
\end{solution}

\begin{problem}
集合 \(\{1,2,3\}\) 的子集共有（\quad）

\choices
    {\(5\) 个}
    {\(6\) 个}
    {\(7\) 个}
    {\(8\) 个}
\end{problem}

\begin{answer}
D
\end{answer}

\begin{solution}
构造一个子集时，元素 \(1\)、\(2\)、\(3\) 分别都有“取入”和“不取入”两种相互独立的选择. 因而子集与这 \(3\) 个元素的取舍方案一一对应，子集总数为 \(2\times2\times2=2^3=8\)，其中包括空集和原集合，故选 D.
\end{solution}

\begin{problem}
函数 \(y = a^x\)（\(0 < a < 1\)）的图象是（\quad）

\choices
    {\choicebitmap[width=0.15\paperwidth]{img/1988/national_liberal/q4_optA.png}}
    {\choicebitmap[width=0.15\paperwidth]{img/1988/national_liberal/q4_optB.png}}
    {\choicebitmap[width=0.15\paperwidth]{img/1988/national_liberal/q4_optC.png}}
    {\choicebitmap[width=0.15\paperwidth]{img/1988/national_liberal/q4_optD.png}}
\end{problem}

\begin{answer}
B
\end{answer}

\begin{solution}
由于 \(0<a<1\)，当 \(x_1<x_2\) 时，\(a^{x_2-x_1}<1\)，从而
\(a^{x_2}=a^{x_1}a^{x_2-x_1}<a^{x_1}\)，所以函数在 \(\R\) 上单调递减. 另外，\(a^x>0\)，且 \(a^0=1\)；当 \(x\to+\infty\) 时，\(a^x\to0\)，当 \(x\to-\infty\) 时，\(a^x\to+\infty\). 只有第二个图象同时满足这些性质，因此选 B.
\end{solution}

\begin{problem}
已知椭圆方程 \(\frac{x^2}{20} + \frac{y^2}{11} = 1\)，那么它的焦距是（\quad）

\choices
    {\(6\)}
    {\(3\)}
    {\(2\sqrt{31}\)}
    {\(\sqrt{31}\)}
\end{problem}

\begin{answer}
A
\end{answer}

\begin{solution}
方程中分母较大的 \(20\) 对应长半轴，所以 \(a^2=20\)，而 \(b^2=11\). 椭圆的焦半距满足
\(c^2=a^2-b^2=20-11=9\)，由于 \(c>0\)，得 \(c=3\). 两焦点间的距离是焦距 \(2c=6\)，故选 A.
\end{solution}

\begin{problem}
在复平面内，与复数 \(z = -1 - \i\) 的共轭复数对应的点位于（\quad）

\choices
    {第一象限}
    {第二象限}
    {第三象限}
    {第四象限}
\end{problem}

\begin{answer}
B
\end{answer}

\begin{solution}
共轭运算只改变虚部的符号，所以
\(\overline z=-1+\i\). 在复平面中，该复数对应点的坐标为 \((-1,1)\)，其横坐标为负、纵坐标为正，因此点位于第二象限，故选 B.
\end{solution}

\begin{problem}
在 \(\left(x-\sqrt{3}\right)^{10}\) 的展开式中，\(x^6\) 的系数是（\quad）

\choices
    {\(-27\mathrm C_{10}^6\)}
    {\(27\mathrm C_{10}^4\)}
    {\(-9\mathrm C_{10}^6\)}
    {\(9\mathrm C_{10}^4\)}
\end{problem}

\begin{answer}
D
\end{answer}

\begin{solution}
按二项式定理，展开式的第 \(k+1\) 项为
\(\mathrm C_{10}^k x^{10-k}\left(-\sqrt{3}\right)^k\)，其中 \(k=0\)，\(1\)，\(\ldots\)，\(10\). 要得到含 \(x^6\) 的项，必须满足 \(10-k=6\)，所以 \(k=4\). 因而该项为
\(\mathrm C_{10}^4x^6\left(-\sqrt{3}\right)^4=9\mathrm C_{10}^4x^6\)，故 \(x^6\) 的系数为 \(9\mathrm C_{10}^4\)，选 D.
\end{solution}

\begin{problem}
函数 \(y = 3\cos\left(\frac{2}{5}x-\frac{\pi}{6}\right)\) 的最小正周期是（\quad）

\choices
    {\(\frac{2}{5}\pi\)}
    {\(\frac{5}{2}\pi\)}
    {\(2\pi\)}
    {\(5\pi\)}
\end{problem}

\begin{answer}
D
\end{answer}

\begin{solution}
余弦函数本身的最小正周期是 \(2\pi\). 令自变量为
\(u=\dfrac25x-\dfrac\pi6\)，当 \(x\) 增加 \(T\) 时，要使函数值第一次重复，应有 \(\dfrac25T=2\pi\). 因此
\(T=\dfrac{2\pi}{\frac25}=5\pi\)，故选 D.
\end{solution}

\begin{problem}
\(\sin\left(-\frac{19}{6}\pi\right)\) 的值等于（\quad）

\choices
    {\(\frac{1}{2}\)}
    {\(-\frac{1}{2}\)}
    {\(\frac{\sqrt{3}}{2}\)}
    {\(-\frac{\sqrt{3}}{2}\)}
\end{problem}

\begin{answer}
A
\end{answer}

\begin{solution}
正弦函数的周期是 \(2\pi\)，而
\(-\dfrac{19\pi}{6}+4\pi=\dfrac{5\pi}{6}\). 所以
\(\sin\left(-\dfrac{19\pi}{6}\right)=\sin\left(\dfrac{5\pi}{6}\right)\). 角 \(\dfrac{5\pi}{6}\) 位于第二象限，参考角为 \(\dfrac\pi6\)，故
\(\sin\left(\dfrac{5\pi}{6}\right)=\sin\dfrac\pi6=\dfrac12\)，选 A.
\end{solution}

\begin{problem}
直线 \(x+ay = 2a+2\) 与 \(ax+y = a+1\) 平行（不重合）的充要条件是（\quad）

\choices
    {\(a = \frac{1}{2}\)}
    {\(a = -\frac{1}{2}\)}
    {\(a = 1\)}
    {\(a = -1\)}
\end{problem}

\begin{answer}
C
\end{answer}

\begin{solution}
两条直线的法向量分别为 \(\left(1,a\right)\) 和 \(\left(a,1\right)\). 它们平行的必要条件是法向量平行，即
\(\begin{vmatrix}1&a\\a&1\end{vmatrix}=1-a^2=0\)，所以只能有 \(a=1\) 或 \(a=-1\).

当 \(a=1\) 时，两条直线分别为 \(x+y=4\) 和 \(x+y=2\)，左端系数相同而常数项不同，因此平行且不重合. 当 \(a=-1\) 时，第一条直线为 \(x-y=0\)，第二条直线为 \(-x+y=0\)，二者表示同一条直线，属于重合情形. 因而满足“平行且不重合”的充要条件是 \(a=1\)，选 C.
\end{solution}

\begin{problem}
函数 \(y = \frac{x-2}{2x-1}\)（\(x \in \R\)，\(x \neq \frac{1}{2}\)）的反函数是（\quad）

\choices
    {\(y = \frac{x-2}{2x-1}\)（\(x \in \R\)，\(x \neq \frac{1}{2}\)）}
    {\(y = \frac{2x-1}{x-2}\)（\(x \in \R\)，\(x \neq 2\)）}
    {\(y = \frac{x+2}{2x-1}\)（\(x \in \R\)，\(x \neq \frac{1}{2}\)）}
    {\(y = \frac{2x-1}{x+2}\)（\(x \in \R\)，\(x \neq -2\)）}
\end{problem}

\begin{answer}
A
\end{answer}

\begin{solution}
设原函数为 \(y=\dfrac{x-2}{2x-1}\)，其中 \(x\ne\dfrac12\). 交叉相乘得
\(y\left(2x-1\right)=x-2\)，整理为 \(x\left(2y-1\right)=y-2\). 若 \(y=\dfrac12\)，左边为零而右边为 \(-\dfrac32\)，方程无解，所以原函数的值域不含 \(\dfrac12\). 对于 \(y\ne\dfrac12\)，有
\(x=\dfrac{y-2}{2y-1}\). 交换自变量与函数值的字母，得到反函数
\(y=\dfrac{x-2}{2x-1}\)，其定义域为 \(x\in\R\) 且 \(x\ne\dfrac12\)，故选 A.
\end{solution}

\begin{problem}
如图，正四棱台中，\(A'D'\) 所在的直线与 \(BB'\) 所在的直线是（\quad）

\bitmapfigure[max width=0.42\linewidth]{img/1988/national_liberal/q12_fig1.png}

\choices
    {相交直线}
    {平行直线}
    {不互相垂直的异面直线}
    {互相垂直的异面直线}
\end{problem}

\begin{answer}
C
\end{answer}

\begin{solution}
设下底面正方形的中心为坐标原点，底面在平面 \(z=0\) 内，并取
\(A=(-u,-u,0)\)、\(B=(u,-u,0)\)、\(C=(u,u,0)\)、\(D=(-u,u,0)\)，上底面相应顶点为
\(A'=(-v,-v,h)\)、\(B'=(v,-v,h)\)、\(C'=(v,v,h)\)、\(D'=(-v,v,h)\)，其中 \(u>0\)、\(v>0\)、\(h>0\)，且正四棱台的上下底面大小不同，故 \(u\ne v\).

直线 \(A'D'\) 在平面 \(z=h\) 内，方向向量可取 \(\overrightarrow{A'D'}=(0,2v,0)\)，且它所在的直线为 \(x=-v\). 直线 \(BB'\) 与平面 \(z=h\) 的交点是 \(B'\)，而 \(B'\) 的第一坐标为 \(v\)，不满足 \(x=-v\)，所以两直线不相交. 又 \(\overrightarrow{BB'}=(v-u,u-v,h)\) 的第三坐标为 \(h\ne0\)，它不可能与 \(\overrightarrow{A'D'}\) 平行，因此两直线不平行. 不相交且不平行，说明两直线是异面直线.

最后，两个方向向量的数量积为
\(\overrightarrow{A'D'}\cdot\overrightarrow{BB'}=2v(u-v)\ne0\). 所以两直线不垂直，结论为不互相垂直的异面直线，选 C.
\end{solution}

\begin{problem}
函数 \(y = \sin\left(x+\frac{\pi}{4}\right)\) 在闭区间（\quad）

\choices
    {\(\left[-\frac{\pi}{2}, \frac{\pi}{2}\right]\) 上是增函数}
    {\(\left[-\frac{3\pi}{4}, \frac{\pi}{4}\right]\) 上是增函数}
    {\([-\pi,0]\) 上是增函数}
    {\(\left[-\frac{\pi}{4}, \frac{3\pi}{4}\right]\) 上是增函数}
\end{problem}

\begin{answer}
B
\end{answer}

\begin{solution}
令 \(t=x+\dfrac\pi4\)，则 \(y=\sin t\). 正弦函数在
\(\left[-\dfrac\pi2,\dfrac\pi2\right]\) 上单调递增.

对选项 B，若 \(x\in\left[-\dfrac{3\pi}{4},\dfrac\pi4\right]\)，则
\(t\in\left[-\dfrac\pi2,\dfrac\pi2\right]\)，所以 \(y\) 在该区间上单调递增.

其余选项均不能成立：选项 A 对应的 \(t\) 区间为 \(\left[-\dfrac\pi4,\dfrac{3\pi}4\right]\)，跨过 \(\dfrac\pi2\)；选项 C 对应的 \(t\) 区间为 \(\left[-\dfrac{3\pi}4,\dfrac\pi4\right]\)，跨过 \(-\dfrac\pi2\)；选项 D 对应的 \(t\) 区间为 \([0,\pi]\)，跨过 \(\dfrac\pi2\). 这些区间内正弦函数都先增后减或先减后增，不能保持单调递增，故选 B.
\end{solution}

\begin{problem}
假设在 \(200\) 件产品中有 \(3\) 件次品，现在从中任意抽取 \(5\) 件，其中至少有 \(2\) 件次品的抽法有（\quad）

\choices
    {\(\mathrm C_3^2\mathrm C_{197}^3 + \mathrm C_3^3\mathrm C_{197}^2\) 种}
    {\(\mathrm C_3^2\mathrm C_{197}^3\) 种}
    {\(\mathrm C_{200}^5 - \mathrm C_{197}^5\) 种}
    {\(\mathrm C_{200}^5 - \mathrm C_3^1\mathrm C_{197}^4\) 种}
\end{problem}

\begin{answer}
A
\end{answer}

\begin{solution}
200 件产品中有 3 件次品和 197 件正品，而抽取 5 件时次品数不可能超过 3. 因此“至少有 2 件次品”分成恰有 2 件和恰有 3 件两种互斥情形.

恰有 2 件次品时，从 3 件次品中选 2 件，再从 197 件正品中选 3 件，共有 \(\mathrm C_3^2\mathrm C_{197}^3\) 种；恰有 3 件次品时，从 3 件次品中全选，再从正品中选 2 件，共有 \(\mathrm C_3^3\mathrm C_{197}^2\) 种. 两种情形相加，抽法总数为
\(\mathrm C_3^2\mathrm C_{197}^3+\mathrm C_3^3\mathrm C_{197}^2\)，故选 A.
\end{solution}

\begin{problem}
已知二面角 \(\alpha-AB-\beta\) 的平面角是锐角 \(\theta\)，\(\alpha\) 内一点 \(C\) 到 \(\beta\) 的距离为 \(3\)，点 \(C\) 到棱 \(AB\) 的距离为 \(4\)，那么 \(\tan\theta\) 的值等于（\quad）

\choices
    {\(\frac{3}{4}\)}
    {\(\frac{3}{5}\)}
    {\(\frac{3}{7}\sqrt{7}\)}
    {\(\frac{1}{3}\sqrt{7}\)}
\end{problem}

\begin{answer}
C
\end{answer}

\begin{solution}
过 \(C\) 作平面 \(\gamma\perp AB\)，令 \(D=\gamma\cap AB\)，则 \(CD\subset\alpha\)，且由已知 \(CD=4\). 设 \(l=\gamma\cap\beta\)，因为截面平面 \(\gamma\) 垂直于棱 \(AB\)，所以二面角的平面角 \(\theta\) 就是直线 \(CD\) 与直线 \(l\) 的锐角.

在截面 \(\gamma\) 中，点 \(C\) 到直线 \(l\) 的距离等于点 \(C\) 到平面 \(\beta\) 的距离，均为 \(3\). 若从 \(C\) 向 \(l\) 作垂线，得到直角三角形，其中斜边为 \(CD=4\)，所以
\(4\sin\theta=3\)，即 \(\sin\theta=\dfrac34\). 又因 \(\theta\) 是锐角，\(\cos\theta=\sqrt{1-\sin^2\theta}=\sqrt{1-\dfrac9{16}}=\dfrac{\sqrt7}{4}\). 因而
\(\tan\theta=\dfrac{\sin\theta}{\cos\theta}=\dfrac{3}{\sqrt7}=\dfrac{3\sqrt7}{7}\)，故选 C.
\end{solution}

\section{填空题}

\begin{problem}
求复数 \(\sqrt{3} - \i\) 的模和辐角的主值，分别为\fillinblank{}，\fillinblank{}.
\end{problem}

\begin{answer}
\(2\)，\(\frac{11\pi}{6}\)
\end{answer}

\begin{solution}
设 \(z=\sqrt3-\i\). 由模的定义，
\(\lvert z\rvert=\sqrt{\left(\sqrt3\right)^2+(-1)^2}=\sqrt4=2\).

设辐角主值为 \(\varphi\). 因为实部为正、虚部为负，向量位于第四象限；又
\(\cos\varphi=\dfrac{\sqrt3}{2}\)，\(\sin\varphi=-\dfrac12\). 在本题采用的 \(0\leq\varphi<2\pi\) 的主值约定下，第四象限对应的角为 \(\varphi=\dfrac{11\pi}{6}\). 因而所求的模和辐角主值分别为 \(2\)、\(\dfrac{11\pi}{6}\).
\end{solution}

\begin{problem}
解方程：\(9^{-x} - 2 \cdot 3^{1-x} = 27\)，解为\fillinblank{}.
\end{problem}

\begin{answer}
\(-2\)
\end{answer}

\begin{solution}
令 \(t=3^{-x}\)，则 \(t>0\)，并且
\(9^{-x}=\left(3^2\right)^{-x}=3^{-2x}=t^2\)，\(2\cdot3^{1-x}=2\cdot3\cdot3^{-x}=6t\). 原方程化为
\(t^2-6t=27\)，即
\(t^2-6t-27=(t-9)(t+3)=0\). 所以 \(t=9\) 或 \(t=-3\)，由 \(t>0\) 舍去 \(-3\)，得到 \(3^{-x}=9=3^2\)，于是 \(-x=2\)，即 \(x=-2\). 代回原方程，左端为 \(81-54=27\)，满足原方程.
\end{solution}

\begin{problem}
已知 \(\sin\theta = -\frac{3}{5}\)，\(3\pi < \theta < \frac{7\pi}{2}\)，求 \(\tan\frac{\theta}{2}\) 的值为\fillinblank{}.
\end{problem}

\begin{answer}
\(-3\)
\end{answer}

\begin{solution}
由 \(3\pi<\theta<\dfrac{7\pi}{2}\)，减去 \(2\pi\) 后得到
\(\pi<\theta-2\pi<\dfrac{3\pi}{2}\)，所以 \(\theta\) 的终边在第三象限，\(\cos\theta<0\). 因而
\(\cos^2\theta=1-\sin^2\theta=1-\dfrac9{25}=\dfrac{16}{25}\)，结合符号得 \(\cos\theta=-\dfrac45\).

半角公式给出
\(\tan\dfrac\theta2=\dfrac{\sin\theta}{1+\cos\theta}=\dfrac{-\frac35}{1-\frac45}=\dfrac{-\frac35}{\frac15}=-3\)，所以填 \(-3\).
\end{solution}

\begin{problem}
一个直角三角形的两条直角边的长分别为 \(3\mathrm{~cm}\) 和 \(4\mathrm{~cm}\)，将这个直角三角形以斜边为轴旋转一周，求所得旋转体的体积为\fillinblank{}.
\end{problem}

\begin{answer}
\(\frac{48}{5}\pi\mathrm{~cm}^3\)
\end{answer}

\begin{solution}
斜边长为
\(c=\sqrt{3^2+4^2}=5\mathrm{~cm}\). 设直角顶点到斜边的垂足为 \(H\)，则由面积的两种表示方法，
\(\dfrac12\cdot3\cdot4=\dfrac12\cdot5\cdot CH\)，所以旋转半径
\(CH=\dfrac{12}{5}\mathrm{~cm}\).

三角形绕斜边旋转后，以 \(H\) 为公共底面圆心，分别形成以斜边两端为顶点的两个圆锥. 若 \(AH=p\)、\(BH=q\)，则两个圆锥的高之和为 \(p+q=AB=5\mathrm{~cm}\). 因而旋转体体积为
\(V=\dfrac13\pi\left(\dfrac{12}{5}\right)^2p+\dfrac13\pi\left(\dfrac{12}{5}\right)^2q=\dfrac13\pi\left(\dfrac{12}{5}\right)^2\left(p+q\right)=\dfrac{48}{5}\pi\mathrm{~cm}^3\).
\end{solution}

\begin{problem}
求 \(\displaystyle\lim\limits_{n\to\infty}\frac{3n^2 + 2n}{n^2 + 3n - 1}=\)\fillinblank{}.
\end{problem}

\begin{answer}
\(3\)
\end{answer}

\begin{solution}
因为 \(n\to\infty\) 时 \(n\ne0\)，分子、分母同除以 \(n^2\)，得到
\(\displaystyle\lim_{n\to\infty}\frac{3+\frac2n}{1+\frac3n-\frac1{n^2}}\). 其中 \(\dfrac2n\to0\)、\(\dfrac3n\to0\)、\(\dfrac1{n^2}\to0\)，且分母的极限为 \(1\ne0\)，所以
\(\displaystyle\lim_{n\to\infty}\frac{3n^2+2n}{n^2+3n-1}=\frac{3+0}{1+0-0}=3\).
\end{solution}

\section{解答题}

\begin{problem}
证明：\(\cos 3\alpha = 4\cos^3\alpha - 3\cos\alpha\).
\end{problem}

\begin{solution}
由三角函数的和角公式，\(\cos3\alpha=\cos(2\alpha+\alpha)=\cos2\alpha\cos\alpha-\sin2\alpha\sin\alpha\). 再利用二倍角公式 \(\cos2\alpha=2\cos^2\alpha-1\) 和 \(\sin2\alpha=2\sin\alpha\cos\alpha\)，有
\(\cos3\alpha=(2\cos^2\alpha-1)\cos\alpha-2\sin^2\alpha\cos\alpha\).

由 \(\sin^2\alpha=1-\cos^2\alpha\)，继续化简得
\[
\cos3\alpha=(2\cos^2\alpha-1)\cos\alpha-2(1-\cos^2\alpha)\cos\alpha=2\cos^3\alpha-3\cos\alpha+2\cos^3\alpha=4\cos^3\alpha-3\cos\alpha
\]
因此等式成立.
\end{solution}

\begin{problem}
如图，四棱锥 \(S - ABCD\) 的底面是边长为 \(1\) 的正方形，侧棱 \(SB\) 垂直于底面，并且 \(SB = \sqrt{3}\)，用 \(\alpha\) 表示 \(\angle ASD\)，求 \(\sin\alpha\) 的值.

\bitmapfigure[width=5.0cm]{img/1988/national_liberal/q22_fig1.png}
\end{problem}

\begin{solution}
以 \(B\) 为原点，取 \(BA\)、\(BC\)、\(BS\) 分别为 \(x\)、\(y\)、\(z\) 轴的正向. 因为底面是边长为 \(1\) 的正方形，且 \(SB\perp\) 底面、\(SB=\sqrt3\)，所以
\(B=(0,0,0)\)、\(A=(1,0,0)\)、\(C=(0,1,0)\)、\(D=(1,1,0)\)、\(S=(0,0,\sqrt3)\).

角 \(ASD\) 的两条边对应的向量为 \(\overrightarrow{SA}=A-S=(1,0,-\sqrt3)\)、\(\overrightarrow{SD}=D-S=(1,1,-\sqrt3)\). 因此
\(\overrightarrow{SA}\cdot\overrightarrow{SD}=1+3=4\)，并且 \(\lvert\overrightarrow{SA}\rvert=\sqrt{1+3}=2\)，\(\lvert\overrightarrow{SD}\rvert=\sqrt{1+1+3}=\sqrt5\).

由向量夹角公式，\(\cos\alpha=\dfrac{\overrightarrow{SA}\cdot\overrightarrow{SD}}{\lvert\overrightarrow{SA}\rvert\lvert\overrightarrow{SD}\rvert}=\dfrac4{2\sqrt5}=\dfrac{2\sqrt5}{5}\). 因为 \(0<\alpha<\pi\)，\(\sin\alpha>0\)，所以
\(\sin\alpha=\sqrt{1-\cos^2\alpha}=\sqrt{1-\dfrac45}=\dfrac{\sqrt5}{5}\).
\end{solution}

\begin{problem}
在双曲线 \(x^2 - y^2 = 1\) 的右支上求点 \(P(a,b)\)，使该点到直线 \(y = x\) 的距离为 \(\sqrt{2}\).
\end{problem}

\begin{solution}
直线 \(y=x\) 可写为 \(x-y=0\)，所以点 \(P(a,b)\) 到该直线的距离为 \(\dfrac{\lvert a-b\rvert}{\sqrt{1^2+(-1)^2}}=\dfrac{\lvert a-b\rvert}{\sqrt2}\). 由题意 \(\dfrac{\lvert a-b\rvert}{\sqrt2}=\sqrt2\)，从而 \(\lvert a-b\rvert=2\)，需要分两种情况.

当 \(a-b=2\) 时，双曲线方程给出 \(1=a^2-b^2=(a-b)(a+b)=2(a+b)\)，所以 \(a+b=\dfrac12\). 联立两式，
\(a=\dfrac{(a-b)+(a+b)}2=\dfrac{2+\frac12}{2}=\dfrac54\)，\(b=\dfrac{(a+b)-(a-b)}2=\dfrac{\frac12-2}{2}=-\dfrac34\). 此时 \(a=\dfrac54>0\)，确实在双曲线右支上.

当 \(a-b=-2\) 时，\(1=(a-b)(a+b)=-2(a+b)\)，所以 \(a+b=-\dfrac12\). 于是 \(a=\dfrac{-2-\frac12}{2}=-\dfrac54<0\)，该点在左支上，不符合右支条件.

因此所求点为 \(P\left(\dfrac54,-\dfrac34\right)\).
\end{solution}

\begin{problem}
解不等式：\(\lg\left(x-\frac{1}{x}\right) < 0\).
\end{problem}

\begin{solution}
因为 \(\lg u\)（底数为 \(10\)）在 \(u>0\) 上单调递增，所以原不等式等价于 \(0<x-\dfrac1x<1\). 先解定义域条件：\(x-\dfrac1x=\dfrac{(x-1)(x+1)}x>0\). 由关键点 \(-1\)、\(0\)、\(1\) 分段作符号判断，得到 \(x\in(-1,0)\cup(1,+\infty)\).

在区间 \((-1,0)\) 内，解 \(x-\dfrac1x<1\) 时乘以负数 \(x\)，不等号方向改变，得到 \(x^2-x-1>0\). 在区间 \((1,+\infty)\) 内乘以正数 \(x\)，得到 \(x^2-x-1<0\). 方程 \(x^2-x-1=0\) 的根为 \(r_1=\dfrac{1-\sqrt5}{2}\)、\(r_2=\dfrac{1+\sqrt5}{2}\)，且 \(r_1\in(-1,0)\)、\(r_2>1\).

因此在 \((-1,0)\) 内取 \((-1,r_1)\)，在 \((1,+\infty)\) 内取 \((1,r_2)\). 原不等式的解集为 \(\left(-1,\dfrac{1-\sqrt5}{2}\right)\cup\left(1,\dfrac{1+\sqrt5}{2}\right)\).
\end{solution}

\begin{problem}
一个数列 \(\{a_n\}\)：当 \(n\) 为奇数时，\(a_n = 5n + 1\)；当 \(n\) 为偶数时，\(a_n = 2^{\frac n2}\). 求这个数列的前 \(2m\) 项的和（\(m\) 是正整数）.
\end{problem}

\begin{solution}
在前 \(2m\) 项中，奇数项和偶数项各有 \(m\) 项. 对 \(k=1\)，\(2\)，\(\ldots\)，\(m\)，有 \(a_{2k-1}=5(2k-1)+1=10k-4\)，\(a_{2k}=2^{\frac{2k}{2}}=2^k\).

所以 \(\displaystyle S_{2m}=\sum_{n=1}^{2m}a_n=\sum_{k=1}^m\left(a_{2k-1}+a_{2k}\right)=\sum_{k=1}^m(10k-4)+\sum_{k=1}^m2^k\). 其中 \(\displaystyle\sum_{k=1}^m(10k-4)=10\cdot\frac{m(m+1)}2-4m=5m^2+m\)，而 \(\displaystyle\sum_{k=1}^m2^k=\frac{2(2^m-1)}{2-1}=2^{m+1}-2\). 因而 \(\displaystyle S_{2m}=5m^2+m+2^{m+1}-2\).
\end{solution}
